\documentclass[a4paper]{article}

\usepackage[dutch]{babel}
\usepackage{geometry}
\usepackage[utf8]{inputenc}

% dimensions
\geometry{left=3cm, top=3cm, right=3cm, bottom=3cm}

% font
\usepackage{DejaVuSans}
\renewcommand*\familydefault{\sfdefault}
\usepackage[T1]{fontenc}

% image
\graphicspath{ {img/} }

% lists
\providecommand{\tightlist}{%
\setlength{\itemsep}{0pt}\setlength{\parskip}{0pt}}

% preamble
\title{Can't Touch This User Manual}
\author{Tim Visée \& Nathan Bakhuijzen}
\date{October 2018}

\begin{document}

  \pagenumbering{gobble}
  \maketitle
  \begin{figure}[h]
    \centering
    \includegraphics[width=\linewidth]{cant-touch-this}
  \end{figure}
  \clearpage

  \section{Problem Definition}
  \textbf{Controlling a computer or peripheral using touchless controls is
    possible, but what motions are natural, effortless and easy to use?}
  Questions of this nature can be answered using the Can't Touch This platform.
  % TODO: Extend this section

  \subsection{Motivation}
  At the start of this minor, students were given a choice in the subject of the
  research. Mister Hani introduced us to a series of subjects, of which the
  LeapMotion project was the most interesting to us. The idea of the LeapMotion
  was to create or extend existing software to enable people to control a
  computer without touching any peripherals.

  \subsection{Background information}
  The events that lead up to this question came from the medical industry. When
  surgeons perform surgeries, they often have to look through medical files, for
  example, an X-ray scan. This is difficult, as the surgeons most likely has
  dirty hands, and can therefore not touch any peripheral. Instead of viewing
  the documents himself, the surgeon has assistants that help him instead.

  The use of assistants is a solution to the problem, but this requires another
  human, and personnel in the medical world is scarce. We therefore want to
  replace these human assistants with a device that allows surgeons to control
  the computer without touch any peripherals.

  The people that are most interested in the results are software engineers,
  surgeons and hospital managers.

  %%% Goal of the manual
  %%%% Show researchers how to use your platform 
  %%%% Show future developers how to continue your work

  % Installation Guide
  %% Platform and OS
  %% Other software dependencies
  %% External resources needed to run the prototype
  %% External resources needed to develop the prototype
  \section{Installation Guide}

  % User Instructions
  %% Instructions to run the experiment (explained in build plan)
  %% Aimed at non technical researchers
  %% Include screenshots when possible
  %% Bonus points: make a video of the instructions, upload to youtube and
  %% include link in manual

  % Requirements
  %% Here you describe what requirements you have already gathered from the client. 
  %% What requirements are done 
  %% What requirements are still open 
  %% Prioritizing of requirements (MOSCOW)

  % Architecture diagram
  %% Context view
  %% Functional view

  % Domain Model
  %% This is where you give future developers the information they need to
  %% continue the work. 
  %% In the domain model you list all: 
  %% API interfaces 
  %% Subsystems

  % Test Report
  %% Add this section even if you did not do any formal testing on the system.  
  %% What is the quality of the code 
  %% What has already been tested 
  %% What are the known bugs in the system
\end{document}
